\paragraph{距离规律与叠加规律都有条件}
球面面积 $4\pi r^2$ 来自以点声源为中心的同心球面。只有声功率没有被吸收、没有其他声源或反射波叠加，且辐射近似各向同性时，穿过这些球面的总功率才相同，从而得到平方反比律。靠近有限尺寸振源、具有强方向性的声束或封闭房间的混响，均需另行分析。
对相同阻抗中沿同一方向传播的相干平面波，强度可写成
\[
 I=I_1+I_2+2\sqrt{I_1I_2}\cos\phi.
\]
只有相位差在测量时间内充分变化、交叉项平均为零时，才可取 $I=I_1+I_2$。两个等强、不相干声源合成后的声强级增量是 $10\log_{10}2\simeq3.01\,\mathrm{dB}$；这个数来自强度倍增，不是把原来的分贝数再乘2。声压级则用 $20\log_{10}(p_{\rm rms}/p_0)$，指数前的20来自强度与声压平方的关系，使用时还须给出参考声压和介质条件。
